% Anexo obrigatório (fora do limite de páginas): ferramentas, para quê, e um caso em que errou.
\section{Declaração de utilização de assistentes de IA}\label{app:ia}

\paragraph{Ferramentas e finalidade.}
Usámos o assistente Claude para: gerar a primeira versão do código do repositório (carregamento,
divisão, pré-processamento, classificação, tópicos, extração de informação e \emph{notebook} do
BERTimbau); explicar alguns conceitos; redigir um esqueleto deste
artigo; e apoiar a interpretação de alguns resultados, que o grupo reviu e validou. A anotação manual dos 100
documentos para a extração de informação foi feita exclusivamente pelos cinco elementos do grupo. O
assistente esclareceu dúvidas sobre as regras do guia. Todas as decisões foram discutidas pelo
grupo, o código foi executado e verificado nos nossos computadores e o texto foi revisto e reescrito
por nós. Não foram usadas outras ferramentas de inteligência artificial.

\paragraph{Casos em que o assistente errou e como demos por isso.}
\begin{enumerate}
\item \textbf{Versão incompatível de uma biblioteca.} O \texttt{requirements.txt} inicial fixava o
\texttt{scikit-learn} 1.5.2, que não tem pacote pré-compilado para Python 3.14 no Windows. Demos
por isso quando um elemento do grupo tentou instalar o ambiente e a compilação falhou. A correção
(versão 1.7.2) foi validada a confirmar que todas as tabelas davam exatamente os mesmos números.
\item \textbf{Pressuposto errado sobre os dados.} A primeira versão do carregador assumia que a
versão normalizada do corpus tinha exatamente os mesmos ficheiros que a versão completa; o código
falhou em quatro ficheiros com numeração diferente.
\item \textbf{Erros de programação apanhados por testes.} O \emph{notebook} do BERTimbau gravava os
identificadores num formato que o \texttt{numpy} recusa ler, e o gerador de tabelas usava um nome
de coluna (\texttt{pipe}) que coincide com um método do \texttt{pandas}. Ambos foram detetados ao
correr o código (o \emph{notebook} foi testado de ponta a ponta com um modelo minúsculo antes de
gastar GPU). Também um comentário acrescentado pelo assistente ao script do custo computacional
partiu uma linha de código (\texttt{SyntaxError}), o que só detetámos ao correr o script no Windows.
\end{enumerate}
